\documentclass{Academic}

\usepackage{booktabs}
\usepackage{makecell}
\usepackage{graphicx}

\begin{document}
%Easy customisation of title page
%TC:ignore
\myabstract{\small
\noindent\textbf{Resumen}
%Abstract below

La clasificación automática de noticias financieras resulta crucial para anticipar 
movimientos en el mercado. Este estudio evalúa la efectividad de un clasificador 
Naïve Bayes Multinomial para determinar la relevancia financiera de un corpus de 
noticias a partir de su contenido textual. La base de datos se dividió en conjuntos 
de entrenamiento y prueba para validar el ajuste del algoritmo. Los resultados 
revelaron variaciones en el rendimiento según la distribución de los datos, alcanzando 
una exactitud predictiva de 0.579. Estos hallazgos sugieren que, aunque el clasificador 
identifica patrones básicos, la complejidad del vocabulario financiero requiere 
complementar el análisis con técnicas avanzadas de selección de características.}
\renewcommand{\myTitle}{Implementación de clasificador naïve Bayes como predictor del nivel de impacto de noticias financieras}
\renewcommand{\MyAuthors}{%
    Jessica Milagros Cordero Rivera$^{1}$ \quad
    Diego Rodrigo González Ramírez$^{1}$\\[0.2em]
    Jose Ángel Lechuga Carrasco$^{1}$ \quad
    Julio Lizárraga Beltrán$^{1}$
}

\renewcommand{\MyAffiliations}{%
    $^{1}$Escuela de Ingeniería y Ciencias
}

\renewcommand{\AuthorIDs}{%
    A01647083 \;|\; A01647198 \;|\; A00228198 \;|\; A01743628
}

\renewcommand{\ShortAuthors}{Cordero Rivera et al.}

\renewcommand{\Keywords}{}
\maketitle
%\vspace{-1.9em}\noindent\rule{\textwidth}{1pt} %add this line if not using abstract
\onehalfspacing
%TC:endignore

\section{Introducción} 

Los mercados financieros operan alrededor de una gran cantidad continua 
de información, modificando las expectativas de inversionistas y, por lo tanto, 
relacionarse con cambios en precios, volumen de negociación y toma de decisiones 
en general. Dentro de este flujo de información, las noticias financieras constituyen una fuente 
particularmente relevante, la evidencia empírica ha mostrado que el 
contenido lingüístico de las noticias contiene información económicamente 
significativa. Se ha encontrado que un mayor pesimismo en los medios se 
asocia con presión a la baja sobre los precios del mercado y con variaciones 
en el volumen de negociación \cite{tetlock2007media}. Adicional a lo anterior, 
se ha observado que la presencia de lenguaje negativo en noticias podía 
aportar información sobre sus resultados contables y rendimientos 
bursátiles posteriores \cite{tetlock2008more}. Además, se ha documentado 
que la información contenida en las noticias puede conservar 
capacidad predictiva sobre los rendimientos del mercado \cite{heston2017news}.

Sin embargo, el análisis de noticias financieras convella varias complicaciones. 
A diferencia de variables numéricas, como precios o rendimientos, 
los titulares corresponden a una cadena de palabras que pueden contener 
ambigüedad y/o dependencia del contexto. Es bien sabido en la literatura previa 
que la orientación semántica general de una frase puede ser distinta de la 
polaridad previa de las palabras individuales. \cite{malo2014good}. Por ello, 
identificar si una noticia puede generar un impacto relevante en el mercado 
constituye un problema más complejo que identificar únicamente si su lenguaje 
es positivo, negativo o neutral.


El clasificador Naïve Bayes es un algoritmo de aprendizaje automático supervisado que se utiliza para tareas de 
clasificación, como la clasificación de texto. Su funcionamiento parte 
de la estimación de la probabilidad de pertenencia a una clase a partir 
de las características observadas y de una hipótesis simplificadora de 
independencia condicional entre los predictores \cite{kavlakoglu_naivebayes}. %67 XD
En aplicaciones de cadenas de texto, las palabras pueden representarse 
mediante su frecuencia dentro de cada obvservación. Se ha mostrado que el modelo 
multinomial de Naïve Bayes presenta un comportamiento efectivo para tareas 
de clasificación de texto basadas en frecuencias de palabras \cite{mccallum1998comparison}. 
Su simplicidad omputacional, capacidad para trabajar con matrices dispersas y fácil 
interpretación hacen de este algoritmo uno apropiado para estudiar titulares financieros. 
Además, ya se han aplicado modelos Naïve Bayes sobre información de noticias 
financieras y se encontró que ésta permitía predecir con mayor precisión 
la dirección de cambios en la volatilidad que la dirección del precio de 
cierre \cite{atkins2018financial}. Esto muestra la utilidad de este clasificador 
en el contexto de noticias financieras.

Este trabajo busca determinar si la información contenida en los 
titulares de noticias financieras, junto con las variables contextuales 
disponibles, permiten discriminar entre noticias de 
\textit{Impacto Alto} y noticias de \textit{Impacto Medio/Bajo}. 
Se propone como hipótesis que la distribución de determinadas 
palabras en los titulares, junto con las características financieras disponibles 
contiene información suficiente para identificar noticias pertenecientes a la categoría de
impacto, ya sea alto o medio/bajo.

El objetivo de este estudio es desarrollar y evaluar un 
clasificador Naïve Bayes capaz de estimar la categoría de impacto de una 
noticia financiera a partir de la información contenida en su titular y de 
variables seleccionadas. Para ello, los titulares se procesan 
mediante técnicas de Procesamiento del Lenguaje Natural (NLP) que incluyen 
tokenización, eliminación de palabras de uso común y construcción de una matriz 
dispersa de frecuencias de términos. Sobre esta representación se entrena un modelo 
de clasificación y se mide su desempeño mediante \textit{accuracy}, precisión, \textit{recall}, 
F1-score y matriz de confusión.


\section{Metodología}

\subsection{Enfoque y tipo de estudio}
Se empleó un enfoque cuantitativo, correlacional, no experimental, basado en el uso de 
GBN para encontrar las relaciones de dependencia probabilística entre los biomarcadores.

\subsection{Fuente de datos}
Los datos fueron proporcionados por una base de datos que contiene informacion sobre titulares de noticias
financieras entre febrero y agosto de 2025. La base de datos fue proporiconada por nuestro profesor de el
Instituto Tecnologico de Estudios Superiores Monterrey


\subsection{Herramientas}
La realización y evaluación del modelo se realizaron utilizando entorno de desarrollo 
integrado \texttt{RStudio} (v2026.08.2), utilizando el lenguaje de programación 
estadístico \texttt{R} (v4.5.3) \cite{R} mediante el sistema de código abierto 
Quarto (v1.8.25), utilizando los paquetes \texttt{bnlearn} (v5.2.1)\cite{Bnlearn}, 
\texttt{tidyverse} (v2.0.0),  \texttt{tidytext} (v0.4.3),  \texttt{stringr} (v1.6.0),  
\texttt{ggplot2} (v4.0.3) y \texttt{dplyr} (v1.2.1).

\subsection{Preprocesamiento}
Se exploro la distribucion de nuestra variable objetivo (el nivel de impacto) para poder facilitarnos el modelado
inicial o para evitar desbalances severos. Simplificamos el prblema recodificando la variable objetivo a dos categoricas 
(impacto alto e imapcto bajo). Se uilizo la funcion \texttt{unnes_tokens()} para descomponer los tiulares de noticias en
palabras individuales, se removieron las palabras comunes en ingles utilizando el dataset stop words.


\subsection{Construcción y selección del modelo global}
Construimos la matriz de frecuencias de palabras donde cada renglón represente un titular y las columnas
correspondan al vocabulario retenido. Dividimos la matriz de datos en training y test set, despues se
entrenó un clasificador Naïve Bayes multinomial, utilizando el training set, para estimar la probabilidad
de pertenencia de cada titular a la categoría de Impacto Alto o Impacto Medio/Bajo.


\subsection{Estrategia analítica para consultas probabilísticas}
Hicimos pruebas utilizando el conjunto de prueba, calculamos el accuracy, precision, recall, 
F1-score y matriz de confusión para poder identifiar los patrones de acierot y error del modelo.


\section{Aplicación}

\subsection{Resultados}


\subsection{Discuciones}


\section{Conclusiones}
El clasificador Naïve Bayes multinomial entrenado con obtuvo un desempeño algo débil con una accuracy de 57.87\%,
que quedó por debajo de la tasa de no informacion de 68.06\%, o sea que, la que se obtiene con solo predecir siempre
la clase mayoritaria. El Kappa de 0.0457 nos reafirma que el modelo no concuerda mucho con la realidad osea que
es casi igual a hacerlo al azar, y la exactitud balanceada de 52.32\% refuerza esta conclusion.

Para la clase de interes (Impacto Alto) el modelo mostro una sensibilidad de 36.96\%, con una precision de 34.9\% y
un F1 de 35.9\%. De 292 titulares clasificados como Impacto Alto, solo 102 lo eran en realidad y se dejaron pasar 174
casos reales de Impacto Alto.

Esto muestra que la idea de tratar cada palabra como si no tuviera relación con las demas, junto con solo fijarse en
si una palabra aparece o no (sin importar cuántas veces se repite), no le da al modelo suficiente informacion para 
distinguir entre las dos categorias de impacto. En resumen el Naïve Bayes no logro hacerlo mejor que si simplemente
se eligiera siempre la categoria mas comun, asi que la forma en que se representaron las palabras no fue suficiente para
esta ocasion.

%TC:ignore
%\clearpage %add new page for references
\singlespacing
\emergencystretch 3em
\hfuzz 1px
\printbibliography[heading=bibnumbered]

% \clearpage
% \begin{appendices}

% \section{Here go any appendices!}

% \end{appendices}

%TC:endignore
\end{document}
